%! Author = abenjeddou
%! Date = 29/05/2026

% Preamble
Today, with the advancement of the Internet of Things (IoT), streaming applications, and the growing need for real-time data processing, implementing an efficient real-time data processing workflow has become essential. Businesses are seeking to maximize the use of real-time data to make informed decisions, improve their operations, and provide more responsive customer services. It is in this context that this final-year project focuses on the design and implementation of a real-time data processing workflow. Faced with these challenges, Apache Flink has emerged as one of the leading real-time data processing platforms, offering high-level processing capabilities, horizontal scalability, and low latency. This project explores how Flink can be used to meet the specific needs of our use case.
This report is structured around eight chapters: the first chapter is dedicated to presenting the general context of our project, including the host organization, the problem statement, and the proposed solution.

The second chapter presents the requirements analysis, where we specify the functional and non-functional requirements of our project.
The third chapter defines the target solution architecture in terms of Apache Flink framework, Kubernetes deployment platform, and the global system design.
The remaining chapters are dedicated to presenting the work carried out throughout six sprints: project setup and message ingestion, then subscription, business rules and notification delivery grouped together as the data processing aspect of the workflow, followed by the AI prediction module, containerization and deployment, and finally observability and monitoring.
